\section{Preliminaries and Setting}
\label{sec:prelim}

\subsection{Fields and the cell ring}

Fix the binary extension field
\[
  \K \;:=\; \Ftwo[X]\big/\bigl(P(X)\bigr),
  \qquad
  P(X) \;=\; X^{128} + X^7 + X^2 + X + 1,
\]
the field of order $2^{\Kbits}$ defined by the GHASH/AES-GCM reduction
pentanomial. Elements of $\K$ are stored as $128$-bit vectors;
multiplication is a carryless product followed by reduction modulo $P$,
and addition is bitwise XOR. We write $\K$ abstractly as the protocol's
\emph{challenge and evaluation field}; the choice of $\Kbits = 128$ over
a wider field is a performance trade-off (it halves the carryless
multiplication count relative to a $192$-bit field) and enters the
analysis only through the soundness denominator $2^{\Kbits}$.

Fix a cell-width parameter $\Dval = 32$ and define the \emph{cell ring}
\[
  \cellring \;:=\; \Ftwo[X]\big/(X^{\Dval}),
\]
a free $\Ftwo$-module of rank $\Dval$. A cell $w \in \cellring$ is a
\emph{bit-polynomial} $w = \sum_{b=0}^{\Dval-1} w_b X^b$ with
$w_b \in \{0,1\}$; we identify $\cellring$ with $\polyat{\Ftwo}{\Dval}$
throughout. The quotient by $(X^{\Dval})$ is structural --- it bounds
the cell width --- and is not the source of any reduction in this note.

\begin{remark}[The two roles of the variable $X$]
\label{rem:two-X}
The symbol $X$ plays two roles, distinguished by the surrounding ring.
As the \emph{cell variable} it indexes the $\Dval$ bit positions of a
cell and, after the construction lifts coefficients into $\K$, the
$\Dval$ coordinates of a polynomial in $\Kpoly$. As the
\emph{defining variable} of $\K = \Ftwo[X]/(P)$ it is the formal symbol
in which field elements are represented. The construction only ever
substitutes the cell variable by a field element $\alpha \in \K$; it
never substitutes the defining variable. We keep the single letter $X$
to match the codebase and rely on the ring context to disambiguate;
where a polynomial has $\K$-valued coefficients we write it explicitly
as an element of $\Kpoly$.
\end{remark}

\subsection{Trace, columns, and the matrix view}

Fix the number of variables $\nu$ and set $n = 2^{\nu}$. The committed
witness is a tuple of $K$ trace columns
$w_1, \dots, w_K \in \cellring^{\,n}$, each a length-$n$ vector of cells
indexed by the boolean hypercube $\{0,1\}^{\nu}$.

For the opening we arrange each column $w_g$ as a matrix
\[
  M_{w_g} \;\in\; \cellring^{\,n_{\mathrm{rows}} \times n_{\mathrm{cols}}},
  \qquad
  n = n_{\mathrm{rows}} \cdot n_{\mathrm{cols}},
\]
both factors powers of two. Correspondingly the variable index splits as
$\{0,1\}^{\nu} = \{0,1\}^{\nu_0} \times \{0,1\}^{\nu_1}$ with
$\nu = \nu_0 + \nu_1$; we write a hypercube point as $(i,j)$ with
$i \in \{0,1\}^{\nu_0}$ (the row index) and $j \in \{0,1\}^{\nu_1}$ (the
column index). In the codebase $n_{\mathrm{cols}}$ is the code's row
length $\code{row\_len}$ and $n_{\mathrm{rows}} = n / \code{row\_len}$;
the deployed binary trace uses a small fixed $n_{\mathrm{rows}}$ (e.g.
$8$).

\subsection{The evaluation projection $\psia$}

The Zinc\texttt{+} PIOP reduces ring-typed constraints to the binary
field by a transcript-fresh \emph{evaluation projection}. Fix
$\alpha \in \K$ sampled after the witness is committed, and define the
$\Ftwo$-algebra homomorphism
\[
  \psia : \cellring \longrightarrow \K,
  \qquad
  \psia\Big(\textstyle\sum_b w_b X^b\Big) \;=\; \sum_b w_b\,\alpha^b
  \;=\; w(\alpha).
\]
That is, $\psia$ evaluates the cell's bit-polynomial at the field point
$\alpha$. We use the same symbol for its coordinatewise action
$\psia : \cellring^{\,n} \to \K^{\,n}$. The map is $\Ftwo$-linear,
indeed an $\Ftwo$-algebra map; this is the only property of $\psia$ the
opening uses.

\begin{remark}[$\psia$ extends to $\Kpoly$]
\label{rem:psi-extends}
Extend $\psia$ to $\Ftwo$-linear maps with $\K$-coefficients by
declaring $\psia(X^b) = \alpha^b$ and extending $\K$-linearly: for
$p = \sum_b p_b X^b \in \Kpoly$, set $\psia(p) := \sum_b p_b\,\alpha^b$.
This is evaluation of the cell variable at $\alpha$, now applied to a
$\K$-coefficient polynomial; it agrees with the cell-level $\psia$ on
$\cellring \subset \Kpoly$. We write $\alpha$-projection for this map
and note it is the inner product of the coefficient vector of $p$ with
$(\alpha^0, \dots, \alpha^{\Dval-1})$.
\end{remark}

\subsection{Multilinear extensions and the input claim}

For a vector $v \in \K^{\,n}$, its \emph{multilinear extension} over
$\K$ is the unique multilinear $\mle{v} \in \K[Y_1,\dots,Y_\nu]$ with
$\mle{v}(c) = v_c$ for all $c \in \{0,1\}^{\nu}$; explicitly
\[
  \mle{v}(r) \;=\; \sum_{c \in \{0,1\}^{\nu}} \eq_c(r)\,v_c,
  \qquad
  \eq_c(r) \;=\; \prod_{k=1}^{\nu}\bigl(c_k r_k + (1-c_k)(1-r_k)\bigr).
\]
The upstream sumcheck delivers to the opening, for each committed column
$w_g$, an \emph{MLE evaluation claim}
\[
  a_g \;=\; \mle{\psia(w_g)}(r^*) \;\in\; \K,
  \qquad g = 1, \dots, K,
\]
at a common point $r^* \in \K^{\,\nu}$. The verifier holds
$(r^*, \alpha, a_1, \dots, a_K)$ and the commitment; the prover holds the
full witness and the commit-time codeword data. \emph{The task of the
evaluation argument is to certify each $a_g$ against the commitment.}

\subsection{The eq-tensor}
\label{sec:eqtensor}

Under the matrix view the $\eq$-MLE factorises as a tensor product over
the row/column split:
\[
  \eq_{(i,j)}(r^*) \;=\; q_0[i]\cdot q_1[j],
  \qquad
  q_0 \in \K^{\,n_{\mathrm{rows}}},\;\; q_1 \in \K^{\,n_{\mathrm{cols}}},
\]
where $q_0, q_1$ are the boolean-hypercube $\eq$-vectors of the two
halves of $r^*$ (computed by the codebase routine
\code{build\_eq\_x\_r\_vec}, split as in \code{build\_eq\_tensor\_gf128}).
The matrix form of the claim is the bilinear form
\begin{equation}
  \label{eq:claim-bilinear}
  a_g \;=\; q_0^{\top}\,\psia(M_{w_g})\,q_1
  \;=\; \sum_{i}\sum_{j} q_0[i]\,q_1[j]\,\psia\!\bigl(M_{w_g}[i,j]\bigr)
  \;\in\;\K,
\end{equation}
where $\psia(M_{w_g})$ is the entrywise $\alpha$-evaluation of the cell
matrix. Both parties compute $(q_0, q_1)$ identically from $r^*$; the
eq-tensor is kept in $\K$ (no lift), which is exactly what makes the
bit-sliced opening possible.
